% Cote 119: offprints, reports, courses, CV and letters (typed and printed part): archive cote 119, site page 169 (PDF page 170).
% typescript. Draft compared: none.
% AI-transcribed by Claude Opus 5.5 (claude-opus-5-5, adaptive thinking, effort high), Message Batches API, batch r119_full (msgbatch_01DwZNm5Y7YsKBWyAN4sFs8w), 2026-11-13; not yet checked by eye.
\scanpage{119/169}
\mspage{- 8 -}

4. \uline{Fonctions holomorphes}

Les fonctions dérivables rencontrés en 4.5. p;19 du fascicule
SERIES ENTIERES vont porter désormais le nom plus glorieux de fonctions
holomorphes et le but annoncé dans l'introduction est de montrer que les
fonctions holomorphes sont analytiques ce qui justifiera a posteriori la
terminologie signalée dans le fascicule FONCTIONS ANALYTIQUES p. 3 .
Par abus de notation on notera souvent par la même lettre $f$ une fonction
d'un ouvert $U$ du plan complexe dans le plan complexe et la fonction induite
en ``oubliant'' la structure complexe de U. Soit $Z = X+iY$ . On écrira :
$f(Z) = f(X,Y)$. On désignera de même par $U$ l'ouvert de $R^{2}$ image de $U$ par
l'application canonique $x+iy \rightarrow (x,y)$ .

\uline{Définition 4.1.}

On dit que $f$ est $\mathbb{R}$-différentiable au point $(x_{o},y_{o}) \in U$ s'il
existe une fonction linéaire $(h,k) \rightarrow ak + bk$ de $R^{2}$ dans $\mathbb{C}$ telle que l'on
ait la relation :
\[ (4.1) \quad f(x_{o}+h,\ y_{o}+k) - f(x_{o},y_{o}) = ah + bk + o(\sqrt{h^{2}+k^{2}}) \quad (\text{où } O\ (\sqrt{h^{2}+k^{2}}) \]
désigne un infiniment petit par rapport à $\sqrt{h^{2}+k^{2}}$) . Si $f$ est $\mathbb{R}$-diffé-%
rentiable au point $(x_{o},y_{o})$ les constantes $a$ et $b$ sont déterminées de manière
unique et on a :
\[ a = \frac{\partial f}{\partial x}\,(x_{o},y_{o}) \qquad b = \frac{\partial f}{\partial y}\,(x_{o},y_{o}) \]
Rappelons que l'existence des dérivées partielles au point $(x_{o},y_{o})$ n'est pas
suffisante pour que la fonction soit dérivable en ce point. Mais on \inserted[typed]{a} la ``réci-%
proque'' suivante :
Si $f$ possède des dérivées partielles dans un voisin\unsure{a}ge de $(x_{o},y_{o})$ et si ces
dérivées partielles sont continues au point $(x_{o},y_{o})$ alors $f$ est différentia-%
ble en ce point.

\margin{169}
